% =====================================================================
%  Resume bullet audit: LangChain / LangGraph transcript RAG assistant
%  Compile:  pdflatex rag_project_audit.tex   (run twice for the ToC)
% =====================================================================
\documentclass[11pt,a4paper]{article}

\usepackage[margin=2.2cm]{geometry}
\usepackage{booktabs}
\usepackage{longtable}
\usepackage{array}
\usepackage{xcolor}
\usepackage{enumitem}
\usepackage{listings}
\usepackage[hidelinks]{hyperref}
\usepackage{titlesec}

\definecolor{supported}{RGB}{0,120,60}
\definecolor{partial}{RGB}{190,120,0}
\definecolor{unsupported}{RGB}{180,30,30}
\definecolor{codebg}{RGB}{245,245,245}

\newcommand{\SUP}{\textcolor{supported}{\textbf{SUPPORTED}}}
\newcommand{\PART}{\textcolor{partial}{\textbf{PARTIALLY SUPPORTED}}}
\newcommand{\NOT}{\textcolor{unsupported}{\textbf{NOT SUPPORTED}}}
\newcommand{\UNV}{\textcolor{partial}{\textbf{NOT VERIFIABLE HERE}}}

\lstdefinestyle{code}{
  basicstyle=\ttfamily\footnotesize,
  backgroundcolor=\color{codebg},
  breaklines=true,
  frame=single,
  framesep=4pt,
  rulecolor=\color{gray!40},
  columns=fullflexible,
  keepspaces=true,
  showstringspaces=false,
  xleftmargin=0pt,
}
\lstset{style=code}

\titleformat{\section}{\large\bfseries}{\thesection}{0.7em}{}
\setlist{itemsep=2pt,parsep=0pt,topsep=3pt}

\title{\textbf{Resume Bullet Audit and Evaluation Report}\\[2pt]
\large LangChain / LangGraph Transcript RAG Assistant}
\author{Parikshit M. \quad (M.Tech Research, CDS, IISc Bengaluru)}
\date{22 September 2026}

\begin{document}
\maketitle

\begin{abstract}
\noindent
This document audits the resume bullets for the project \emph{``Retrieval-Augmented
Generation: From-Scratch Implementation and Deployed Document Assistant''} against
the code that actually exists in the repository
\texttt{DL/Lang\_Chain/Lang\_Chain}. Every factual claim below is tied to a file and
line number. Section~\ref{sec:eval} reports retrieval metrics measured by running a
purpose-built evaluation harness on the transcript corpus that this pipeline has
genuinely indexed. Two words on the resume --- \emph{deployed} and \emph{agent} ---
do not survive the audit and are addressed explicitly in
Sections~\ref{sec:deployed} and \ref{sec:agent}.
\end{abstract}

\tableofcontents
\vspace{1em}

% =====================================================================
\section{Scope and method}
% =====================================================================

\paragraph{What was audited.} The two RAG applications in the repository:
\begin{itemize}
  \item \texttt{RAG\_Langgraph/} --- earlier version, in-memory vector store, in-memory checkpointer.
  \item \texttt{Transcript\_RAG\_With\_DB/} --- current version, persistent Chroma
        store plus SQLite checkpointing. \textbf{This is the version the resume
        should describe}; all line references below are to this version unless
        stated otherwise.
\end{itemize}

\paragraph{What could not be audited.} The from-scratch NumPy implementation
referenced in bullet~1 is not present in this repository. Claims about it are
marked \UNV{} rather than unsupported. Section~\ref{sec:eval} ships a harness
written so the \emph{identical} question set and metrics can be run against that
implementation with one command.

\paragraph{What was measured versus estimated.} Retrieval metrics and retrieval
latency in Section~\ref{sec:eval} were measured by execution on the real indexed
corpus. End-to-end latency and answer faithfulness were \emph{not} measured,
because they require live Groq API calls; the API keys in the repository's
\texttt{.env} are exposed (see Section~\ref{sec:defects}) and must be rotated
before any further calls are made. The harness supports those measurements once
new keys are in place.

% =====================================================================
\section{Factual audit of the code}
% =====================================================================

\subsection{Models, store, and chunking}

\begin{longtable}{@{}p{3.6cm}p{7.3cm}p{4.6cm}@{}}
\toprule
\textbf{Question} & \textbf{Answer from the code} & \textbf{Evidence} \\
\midrule
\endhead
Embedding model &
\texttt{all-MiniLM-L6-v2} via \texttt{HuggingFaceEmbeddings}
(sentence-transformers). \textbf{384 dimensions}, max sequence length
\textbf{256 tokens}, outputs are L2-normalised (the model ships a
\texttt{Normalize} module). Runs locally on CPU/GPU; no API call. &
\texttt{Transcript\_RAG\_backend.py:22}; dimension confirmed against the live
Chroma database (\texttt{collections.dimension = 384}). \\
\addlinespace
Chat LLM &
\texttt{openai/gpt-oss-120b} served by \textbf{Groq}, i.e.\ a remote API, not
local, and not OpenAI's own service. The model is OpenAI's open-weight release
hosted on Groq's inference platform. &
\texttt{Transcript\_RAG\_backend.py:21}\newline
\texttt{llm = ChatGroq(model="openai/gpt-oss-120b")} \\
\addlinespace
Vector store &
\textbf{Chroma} (not FAISS). Persisted to disk at
\texttt{Transcript\_RAG\_With\_DB/storage/chroma\_store}, one collection per
transcript. The earlier \texttt{RAG\_Langgraph} version used an ephemeral
in-memory Chroma. &
\texttt{Transcript\_RAG\_backend.py:28,184--191};\newline
in-memory variant at \texttt{RAG\_Langgraph/rag\_backend.py:58} \\
\addlinespace
Similarity metric &
Chroma's default \texttt{l2} space. Because MiniLM vectors are unit-normalised,
L2 ranking is order-equivalent to cosine ranking
($\lVert a-b\rVert^2 = 2-2\cos\theta$). Retrieval is approximate (HNSW), though
at 65 vectors it is effectively exact. &
No \texttt{collection\_metadata} rows set in
\texttt{chroma.sqlite3}, so the default applies; confirmed empirically in
Section~\ref{sec:eval}, where Chroma and exact cosine search return identical
rankings on all 35 questions. \\
\addlinespace
Chunking (\texttt{.txt}, \texttt{.pdf}) &
\texttt{RecursiveCharacterTextSplitter}, \textbf{chunk\_size = 1000 characters,
chunk\_overlap = 200 characters}. &
\texttt{Transcript\_RAG\_Loader.py:8} \\
\addlinespace
Chunking (\texttt{.srt}, \texttt{.vtt}) &
A \textbf{different, hand-written} path: a regex extracts timestamped cues, then
consecutive cues are accumulated until $\geq 1000$ characters. \textbf{No
overlap}, and each chunk carries the timestamp of its first cue in metadata. &
\texttt{Transcript\_RAG\_Loader.py:12--55} \\
\addlinespace
Retrieval depth &
\texttt{k = 4}, fixed, no re-ranking, no MMR, no score threshold. &
\texttt{Transcript\_RAG\_backend.py:231} \\
\addlinespace
Upload formats &
\texttt{.txt}, \texttt{.pdf}, \texttt{.srt}, \texttt{.vtt} --- four formats, all
genuinely handled. \texttt{.csv} and \texttt{.docx} are \emph{not} supported by
this app. &
\texttt{Transcript\_RAG\_Frontend.py:96--99};\newline
\texttt{Transcript\_RAG\_Loader.py:72--95} \\
\bottomrule
\end{longtable}

\paragraph{A chunking detail worth knowing before an interview.}
The chunk size (1000 characters) sits close to the embedding model's hard limit
of 256 tokens. Measured on the indexed corpus: mean 235 tokens per chunk, maximum
263, and \textbf{4 of 65 chunks (6\%) exceed 256 tokens and are silently
truncated} before embedding. On this transcript the loss is trivial (about 2
tokens each), but on denser prose --- a contract or a paper, where characters per
token is lower --- a 1000-character chunk would routinely overflow and the tail of
every chunk would never be embedded. This is a real, defensible observation about
your own system, and naming it first is much stronger than being shown it.

\subsection{LangGraph topology: every node and edge}

\begin{lstlisting}
START  -->  retrieve  -->  generate  -->  END
\end{lstlisting}

\begin{itemize}
  \item \textbf{Nodes (2):} \texttt{retrieve}, \texttt{generate}.
        (\texttt{Transcript\_RAG\_backend.py:248--249})
  \item \textbf{Edges (3):} \texttt{START}$\to$\texttt{retrieve},
        \texttt{retrieve}$\to$\texttt{generate},
        \texttt{generate}$\to$\texttt{END}.
        (\texttt{Transcript\_RAG\_backend.py:250--252})
  \item \textbf{Conditional routing:} none. There is no
        \texttt{add\_conditional\_edges} call anywhere in the repository's RAG code.
  \item \textbf{Query rewriting:} none.
  \item \textbf{Retrieval grading / relevance filtering:} none. Whatever the
        retriever returns is passed to the LLM verbatim.
  \item \textbf{Retry:} none \emph{inside the graph}. There is a retry
        (\texttt{\_invoke\_with\_backoff}, lines 95--107), but it is an ordinary
        Python loop around a Groq call that handles \texttt{RateLimitError}. It is
        not a graph edge and it does not re-retrieve.
  \item \textbf{State:} \texttt{RAGState} carries \texttt{messages} (reduced by
        \texttt{add\_messages}) and \texttt{context}. (lines 59--61)
  \item \textbf{Persistence:} compiled with a \texttt{SqliteSaver} checkpointer,
        so multi-turn history survives process restarts; threads are named
        \texttt{\{transcript\_id\}::\{chat\_id\}}. (lines 31--32, 254;
        \texttt{Transcript\_RAG\_Frontend.py:177})
\end{itemize}

\noindent
\textbf{Verdict: it is a fixed \texttt{retrieve}$\to$\texttt{generate} chain.}
What LangGraph is genuinely buying here is the checkpointer and
node-filtered token streaming --- not branching.

\subsection{Is the answer grounded?}

Yes, by prompt instruction only. The exact constraint, at
\texttt{Transcript\_RAG\_backend.py:49--56}:

\begin{lstlisting}
RAG_PROMPT = ChatPromptTemplate.from_messages([
    ("system",
     "You are a helpful assistant answering questions about a meeting/video transcript. "
     "Use ONLY the transcript excerpts below to answer. If the answer isn't in the "
     "excerpts, say you don't know based on the transcript.\n\n"
     "Transcript excerpts:\n{context}"),
    ("placeholder", "{messages}"),
])
\end{lstlisting}

\noindent What is \emph{not} there:
\begin{itemize}
  \item \textbf{No low-score fallback.} Retrieval has no distance threshold; four
        chunks are always returned and always injected, however irrelevant.
        (line 231)
  \item \textbf{No source display in the UI.} Each excerpt is prefixed with its
        timestamp or filename \emph{inside the prompt}
        (\texttt{Transcript\_RAG\_backend.py:239--242}), so the model \emph{may}
        cite it in prose, but the Streamlit interface renders only the answer
        text. There is no citation panel and no retrieved-chunk display.
        (\texttt{Transcript\_RAG\_Frontend.py:183--193})
  \item \textbf{No faithfulness measurement.} Nothing verifies that the model
        obeyed the instruction.
\end{itemize}

\noindent
Accurate phrasing for the resume: answers are \emph{context-constrained by
prompt}, with timestamp provenance passed to the model. They are not
\emph{citation-verified}.

\subsection{Deployment artefacts}

\begin{longtable}{@{}p{5cm}p{10.5cm}@{}}
\toprule
\textbf{Artefact} & \textbf{Status} \\
\midrule
\endhead
Dockerfile & \textbf{Absent.} No Dockerfile anywhere in the repository. \\
\texttt{requirements.txt} & \textbf{Present but stale and wrong.} The file is named
\texttt{requirement.txt} (singular) and \textbf{omits every dependency this app
actually needs}: \texttt{streamlit}, \texttt{langgraph},
\texttt{langchain-groq}, \texttt{chromadb}, \texttt{sentence-transformers},
\texttt{langgraph-checkpoint-sqlite}. Nothing is version-pinned. The app cannot be
installed from it. \\
Host configuration & \textbf{Absent.} No \texttt{.streamlit/config.toml}, no
\texttt{Procfile}, no \texttt{space.yaml}, no CI workflow. \\
Secrets handling & \texttt{load\_dotenv()} reading a local \texttt{.env}
(\texttt{Transcript\_RAG\_backend.py:17,19}). The \texttt{.env} file is
\textbf{committed to git} and contains five live API keys. See
Section~\ref{sec:defects}. \\
Public URL & \textbf{None.} The app runs only via \texttt{streamlit run} on
localhost. \\
\bottomrule
\end{longtable}

% =====================================================================
\section{Claim-by-claim verdicts}
% =====================================================================

\subsection{Bullet 1 --- from-scratch NumPy implementation}

\begin{quote}\itshape
Built a Retrieval-Augmented Generation (RAG) pipeline from scratch in NumPy,
implementing document chunking, embedding, vector similarity search and context
assembly directly, to establish the underlying mechanics without framework
abstraction.
\end{quote}

\begin{longtable}{@{}p{7.2cm}p{3.2cm}p{5cm}@{}}
\toprule
\textbf{Claim} & \textbf{Verdict} & \textbf{Note} \\
\midrule
\endhead
Pipeline built from scratch in NumPy & \UNV & Code is outside this repository. \\
Chunking implemented directly & \UNV & --- \\
Embedding implemented directly & \UNV & Clarify for yourself whether you wrote the
\emph{embedding model} or called a pretrained encoder and wrote the surrounding
pipeline. The second is normal and defensible; the first is a much larger claim.
An interviewer will ask. \\
Vector similarity search implemented directly & \UNV & --- \\
Context assembly implemented directly & \UNV & --- \\
\bottomrule
\end{longtable}

\noindent
\textbf{Action required:} this bullet is unverifiable from the repository under
review. Before the interview, put that code in a repository you can show, and run
the harness in Section~\ref{sec:eval} against it so bullet~3 becomes real.

\subsection{Bullet 2 --- the LangChain / LangGraph assistant}

\begin{quote}\itshape
Developing a working document assistant using LangChain and LangGraph with a
Streamlit chat interface served locally, supporting document upload, semantic
retrieval through an embedding model, and grounded answer generation through a
chat LLM.
\end{quote}

\begin{longtable}{@{}p{6.4cm}p{3.4cm}p{5.6cm}@{}}
\toprule
\textbf{Claim} & \textbf{Verdict} & \textbf{Evidence} \\
\midrule
\endhead
Uses LangChain & \SUP & \texttt{ChatGroq}, \texttt{HuggingFaceEmbeddings},
\texttt{Chroma}, \texttt{ChatPromptTemplate},
\texttt{RecursiveCharacterTextSplitter}, \texttt{PyPDFLoader}
(\texttt{Transcript\_RAG\_backend.py:8--13}). \\
\addlinespace
Uses LangGraph & \SUP & \texttt{StateGraph} compiled with a checkpointer
(\texttt{Transcript\_RAG\_backend.py:247--254}). \\
\addlinespace
Streamlit chat interface & \SUP & \texttt{st.chat\_input},
\texttt{st.chat\_message}, \texttt{st.write\_stream}
(\texttt{Transcript\_RAG\_Frontend.py:171--193}). \\
\addlinespace
Served locally & \SUP & Accurate, and the honest word. See
Section~\ref{sec:deployed}. \\
\addlinespace
Document upload & \SUP & \texttt{st.file\_uploader} for four formats
(\texttt{Transcript\_RAG\_Frontend.py:96--99}). \\
\addlinespace
Semantic retrieval through an embedding model & \SUP & MiniLM embeddings +
Chroma vector search, $k=4$ (\texttt{Transcript\_RAG\_backend.py:22,231}). \\
\addlinespace
Grounded answer generation & \PART & The prompt constrains the model to the
excerpts (lines 49--56), but there is no retrieval threshold, no citation
rendering, and no faithfulness evaluation. ``Grounded'' overstates a
prompt-level constraint. \\
\addlinespace
``Document assistant'' & \PART & The system prompt is transcript-specific
(``a meeting/video transcript''), the \texttt{.srt}/\texttt{.vtt} path is
purpose-built, and the app auto-generates a meeting summary with action items.
It is a \emph{transcript} assistant. Calling it generic undersells the
specialisation and invites a question you will lose. \\
\bottomrule
\end{longtable}

\noindent
\textbf{Undersold, and worth adding:} the bullet omits your three strongest pieces
of engineering --- SQLite-checkpointed multi-turn memory with resumable named
threads (lines 31--32, 210--226), persistent per-transcript Chroma collections
that avoid re-embedding on re-upload (lines 174--200), and a map-reduce
summarisation path with exponential backoff built specifically to stay under
Groq's tokens-per-minute cap (lines 85--132).

\subsection{Bullet 3 --- the comparison}

\begin{quote}\itshape
Comparing the two implementations on retrieval quality, latency and engineering
effort, and documenting where framework abstractions add value and where they
obscure failure modes.
\end{quote}

\begin{longtable}{@{}p{6.4cm}p{3.4cm}p{5.6cm}@{}}
\toprule
\textbf{Claim} & \textbf{Verdict} & \textbf{Evidence} \\
\midrule
\endhead
Comparing on retrieval quality & \NOT{} (before this report) & No evaluation set,
no metric code, and no results existed anywhere in the repository. \\
\addlinespace
Comparing on latency & \NOT{} (before this report) & No timing instrumentation
existed. \\
\addlinespace
Comparing on engineering effort & \PART & Defensible by line count
(437 non-blank, non-comment lines across three files) but nothing was recorded. \\
\addlinespace
Documenting where abstractions add value / obscure failure modes & \NOT & No such
document existed. This report and Section~\ref{sec:eval} are the first artefact
that substantiates it. \\
\bottomrule
\end{longtable}

\noindent
\textbf{Do not leave bullet 3 as written} unless you run the harness against both
implementations. A claim to have compared systems, with no numbers when asked,
is the worst outcome in an interview.

% =====================================================================
\section{The word ``Deployed''}\label{sec:deployed}
% =====================================================================

\textbf{Ruling: false as the code stands.} There is no Dockerfile, no host
configuration, no usable requirements file and no public URL. The application
runs only under \texttt{streamlit run} on your own machine. Remove the word, or
spend the half-day it takes to make it true.

\subsection{Does the model choice allow a real deployment?}

\textbf{Yes --- this is the good news.} The LLM is a \emph{remote Groq API call},
so no GPU is needed on the host; it needs only an API key and outbound network
access. The embedding model is the constraint: \texttt{all-MiniLM-L6-v2} runs
in-process and pulls in \texttt{torch} plus \texttt{sentence-transformers}
(several hundred megabytes) and roughly 90\,MB of model weights at first run.
Measured cold start for the encoder on this machine: about 35 seconds.

\subsection{What it would actually take}

\begin{longtable}{@{}p{0.6cm}p{6.2cm}p{8.6cm}@{}}
\toprule
& \textbf{Task} & \textbf{Detail} \\
\midrule
\endhead
1 & Rotate all five API keys & Non-negotiable first step. The current keys are in
git history. \\
2 & Write a real \texttt{requirements.txt} & Must add \texttt{streamlit},
\texttt{langgraph}, \texttt{langchain-groq}, \texttt{chromadb},
\texttt{sentence-transformers}, \texttt{langgraph-checkpoint-sqlite},
\texttt{pypdf}, with versions pinned. \\
3 & Move secrets to the platform & \texttt{st.secrets["GROQ\_API\_KEY"]} on
Streamlit Community Cloud, or repository secrets on Hugging Face Spaces, instead
of \texttt{.env}. \\
4 & Solve persistence & \textbf{This is the real obstacle.} Both platforms have
ephemeral filesystems: on redeploy or restart, \texttt{storage/chroma\_store} and
\texttt{chat\_history.db} are wiped, and the ``resume a previous transcript''
feature silently stops working. Options: buy Hugging Face Spaces persistent
storage; or move Chroma to a hosted vector database and the checkpointer to
Postgres (\texttt{langgraph-checkpoint-postgres}). \\
5 & Fix the shared SQLite connection & One module-level connection with
\texttt{check\_same\_thread=False} (line 31) is fine for one local user and a
genuine contention risk with concurrent sessions. \\
6 & Handle multi-user state & \texttt{thread\_id} is derived from the filename
alone, so two visitors uploading files with the same name share a conversation
thread and a vector collection. Add a per-session identifier. \\
7 & Respect the cold start & Cache the encoder with
\texttt{@st.cache\_resource} so it loads once per container, not once per
session. \\
\bottomrule
\end{longtable}

\noindent
\textbf{Platform recommendation: Hugging Face Spaces.} More RAM (16\,GB on the
free CPU tier versus roughly 1\,GB on Streamlit Community Cloud), a model cache
that makes the MiniLM download painless, and purchasable persistent storage,
which is the one thing this app genuinely needs. Streamlit Community Cloud will
struggle with the \texttt{torch} footprint.

% =====================================================================
\section{The word ``Agent''}\label{sec:agent}
% =====================================================================

\textbf{Ruling: not justified. Do not use it.}

An agent chooses what to do next. This graph cannot: it has two nodes joined by
unconditional edges, and every request follows exactly the same path
(\texttt{Transcript\_RAG\_backend.py:247--254}). There is no tool binding, no
\texttt{add\_conditional\_edges}, no loop, no self-correction. The LLM emits text;
it never influences control flow. The only structured decision anywhere is
\texttt{with\_structured\_output} forcing the summary into a \texttt{TypedDict}
(line 79), which is output formatting, not agency.

\noindent
\textbf{Defensible alternatives:} ``stateful LangGraph pipeline'',
``checkpointed retrieve-generate graph'', or simply ``RAG pipeline''. If an
interviewer asks whether it is agentic, the winning answer is: \emph{``No. It is a
deterministic two-node graph. I used LangGraph for the checkpointer and streaming,
not for routing. If I needed agency, the first thing I would add is a retrieval
grading node with a conditional edge back to a query-rewrite node --- the CRAG
pattern.''} That answer demonstrates you know exactly what you did and did not
build.

% =====================================================================
\section{Measured evaluation}\label{sec:eval}
% =====================================================================

\subsection{What was built}

Three files, committed under \texttt{eval/}:

\begin{longtable}{@{}p{4.6cm}p{10.9cm}@{}}
\toprule
\textbf{File} & \textbf{Purpose} \\
\midrule
\endhead
\texttt{corpus\_ai\_infra.json} & The 65 chunks of
\texttt{ai\_infra\_transcript.txt} (61{,}315 characters), extracted from the
Chroma collection your app actually built. The evaluation therefore runs on real
indexed data, not a toy document. \\
\texttt{eval\_questions.json} & 35 questions spanning the whole transcript. \\
\texttt{evaluate\_rag.py} & The harness: pluggable retrieval backends, identical
metrics for each. \\
\bottomrule
\end{longtable}

\paragraph{How gold labels are assigned --- the part an interviewer will probe.}
Ground truth is \emph{not} a hand-typed chunk index. Each question carries an
\texttt{answer\_span}, a short string that states the answer, and every chunk
containing that span is gold. Two consequences matter: labelling is reproducible
and auditable, and the 200-character splitter overlap is handled correctly, since
a fact straddling a boundary legitimately appears in two chunks and both count as
correct. Measured: 1.17 gold chunks per question.

\subsection{Results}

Two backends were run on the same 35 questions over the same 65 chunks, each
query timed three times with the minimum taken:

\begin{itemize}
  \item \textbf{chroma} --- the production path:
        \texttt{HuggingFaceEmbeddings} + \texttt{Chroma} (HNSW), exactly as the
        app builds it.
  \item \textbf{brute} --- \texttt{all-MiniLM-L6-v2} embeddings with exact NumPy
        cosine search, no framework. This is the same shape as your from-scratch
        implementation and acts as the control.
\end{itemize}

\begin{center}
\begin{tabular}{@{}lrrl@{}}
\toprule
\textbf{Metric} & \textbf{Chroma (app)} & \textbf{NumPy (exact)} & \textbf{Note} \\
\midrule
hit@1              & 0.486 & 0.486 & 17 of 35 \\
hit@3              & 0.743 & 0.743 & \\
\textbf{hit@4}     & \textbf{0.829} & \textbf{0.829} & \textbf{29 of 35 --- the app's actual $k$} \\
hit@5              & 0.829 & 0.829 & \\
hit@10             & 0.943 & 0.943 & 33 of 35 \\
\midrule
recall@4           & 0.781 & 0.781 & \\
recall@10          & 0.929 & 0.929 & \\
MRR                & 0.633 & 0.633 & \\
\midrule
Latency mean (CPU) & 9.5\,ms & 7.9\,ms & same run, same machine \\
Latency p95 (CPU)  & 10.7\,ms & 9.3\,ms & \\
Latency mean (GPU) & --- & 3.5\,ms & RTX 6000 Ada \\
\midrule
Encoder cold start & \multicolumn{2}{c}{$\approx$35\,s} & once per process \\
\bottomrule
\end{tabular}
\end{center}

\paragraph{The two backends returned identical rankings on all 35 questions.}
That is a result, not a coincidence, and you should be able to explain it:
MiniLM output is unit-normalised, so Chroma's default L2 ordering is
order-equivalent to cosine, and at 65 vectors HNSW returns the exact neighbours.
The framework therefore buys \emph{zero} retrieval quality at this scale and costs
about 1.6\,ms per query (roughly 20\%) in overhead. It earns its place through
persistence, metadata filtering and the scaling behaviour that would matter at
$10^6$ vectors --- none of which this corpus exercises. That sentence is
exactly what bullet~3 promises to deliver.

\paragraph{The one-line version for an interview.}
\emph{``On 35 questions over a 65-chunk transcript, hit rate at $k=4$ is 0.83 and
MRR is 0.63. Chroma and an exact NumPy search return identical rankings, so the
vector store buys me nothing on quality at this scale and costs about 20\% in
latency; retrieval is under 10\,ms either way, which means essentially all
user-visible latency is the LLM call.''}

\paragraph{Read the failures, not just the score.} Six questions miss at $k=4$,
and the pattern is instructive. For three of them the retriever returns the
\emph{adjacent} chunk --- for instance the question about throughput gains from
splitting prefill and decode pools retrieves chunk 50 when the number is in chunk
51. The discussion is matched correctly; the specific figure lands just over a
boundary. That is a chunking-granularity failure, not a semantic one, and it is
the single strongest argument for the experiment in
Section~\ref{sec:nextsteps}.

\subsection{Running the identical evaluation on the NumPy implementation}

The harness takes the from-scratch pipeline as a plug-in backend. Expose either
of these from a module:

\begin{lstlisting}
# option A
def build_index(chunks: list[str]): ...
def search(index, query: str, k: int) -> list[int]:   # returns chunk indices

# option B
class NumpyRAG:
    def index(self, chunks: list[str]): ...
    def search(self, query: str, k: int) -> list[int]: ...
\end{lstlisting}

\noindent Then run:

\begin{lstlisting}
python eval/evaluate_rag.py --backend brute                      # reference
python eval/evaluate_rag.py --backend chroma                     # the app stack
python eval/evaluate_rag.py --backend numpy --numpy-module numpy_rag
python eval/evaluate_rag.py --backend brute --backend numpy --out results.json
\end{lstlisting}

\noindent
Same corpus, same questions, same gold labels, same metrics --- which is what makes
the comparison in bullet~3 defensible rather than rhetorical.

\subsection{Engineering effort, counted}

\begin{center}
\begin{tabular}{@{}lrr@{}}
\toprule
\textbf{Component} & \textbf{Lines} & \textbf{Non-blank, non-comment} \\
\midrule
\texttt{Transcript\_RAG\_backend.py}  & 254 & 198 \\
\texttt{Transcript\_RAG\_Frontend.py} & 195 & 160 \\
\texttt{Transcript\_RAG\_Loader.py}   &  97 &  79 \\
\midrule
\textbf{Total (framework version)}    & \textbf{546} & \textbf{437} \\
\bottomrule
\end{tabular}
\end{center}

\noindent
Record the same count for the NumPy version. Effort is only a meaningful
comparison if both sides are measured the same way, and line count alone is a
weak proxy --- pair it with what each version does \emph{not} have to implement
(streaming, checkpointing, PDF parsing, persistence).

% =====================================================================
\section{Rewritten bullets}
% =====================================================================

\subsection{Project title}

``Deployed'' does not hold. Options, in order of preference:

\begin{enumerate}
  \item \textbf{Retrieval-Augmented Generation: From-Scratch NumPy Implementation
        and LangGraph Transcript Assistant} \quad --- accurate today.
  \item \textbf{Retrieval-Augmented Generation: From-Scratch and Framework
        Implementations, Benchmarked} \quad --- leads with the comparison, which is
        the most distinctive thing about the project.
  \item Keep ``Deployed'' \emph{only} if you complete the checklist in
        Section~\ref{sec:deployed} and can give a live URL.
\end{enumerate}

\subsection{The three bullets, rewritten}

Each is one line at normal resume font. Numbers are the measured ones; the bracketed
placeholder must be filled by running the harness against your NumPy code.

\paragraph{Bullet 1.}
\begin{quote}
Implemented a RAG pipeline from scratch in NumPy --- chunking, embedding,
vector similarity search and context assembly --- with no retrieval framework, to
expose the mechanics that abstractions hide.
\end{quote}

\paragraph{Bullet 2.}
\begin{quote}
Built a transcript question-answering assistant (LangChain, LangGraph, Streamlit;
MiniLM-384d embeddings, persistent Chroma, Groq \texttt{gpt-oss-120b}) with a
two-node retrieve--generate graph, SQLite-checkpointed multi-turn threads, and
\texttt{.txt}/\texttt{.pdf}/\texttt{.srt}/\texttt{.vtt} ingestion.
\end{quote}

\paragraph{Bullet 3.}
\begin{quote}
Benchmarked both on a 35-question gold set over a 65-chunk corpus: identical
hit@4 (0.83) and MRR (0.63) for Chroma and exact NumPy search, with the framework
costing 20\% retrieval latency [NumPy version: \underline{\quad\quad}], and
identified chunk-boundary misses as the dominant failure mode.
\end{quote}

\paragraph{If you would rather have two bullets than three,} merge 1 and 3: the
comparison is the differentiator, and a from-scratch implementation with no
measurement attached is a weaker claim than a measured one.

% =====================================================================
\section{Ten hard interview questions, with answers from this code}
% =====================================================================

\begin{enumerate}[leftmargin=*,itemsep=6pt]

\item \textbf{Your \texttt{retrieve} node embeds \texttt{state["messages"][-1]}.
What happens when I ask a follow-up like ``and how much memory does that one
have?''}\\
\emph{Answer:} Retrieval degrades, and I know it. Line 235 embeds only the latest
message, so the pronoun is what gets searched; the conversation history reaches the
LLM (via the \texttt{\{messages\}} placeholder, line 55) but never reaches the
retriever. My measured hit@4 of 0.83 is a single-turn number and an upper bound on
multi-turn performance. The standard fix is a condense-question step that rewrites
the follow-up into a standalone query using the history before retrieval.

\item \textbf{Why $k=4$? What did you lose by not choosing 10?}\\
\emph{Answer:} It was a default I then measured. hit@4 is 0.829 and hit@10 is
0.943, so four more of the 35 questions become answerable at $k=10$ --- at the
cost of 2.5$\times$ the context tokens per call, which matters because Groq's free
tier caps me at roughly 8000 tokens per minute. Given that retrieval takes 7\,ms
and the LLM call dominates latency, $k=6$ is probably the better trade, and I have
the harness to confirm it.

\item \textbf{What distance metric does your vector store use? Are you sure?}\\
\emph{Answer:} Chroma's default, which is L2, not cosine --- I set no
\texttt{collection\_metadata}. It does not change the ranking, because
\texttt{all-MiniLM-L6-v2} ships a normalisation layer, so all vectors are unit
length and $\lVert a-b\rVert^2 = 2-2\cos\theta$ is monotonic in cosine
similarity. I verified the vectors are unit-norm rather than assuming it. If I
ever swapped to an unnormalised encoder, the ranking would change and I would need
to set the space explicitly.

\item \textbf{Your chunks are 1000 characters. Your embedding model takes 256
tokens. Walk me through that.}\\
\emph{Answer:} They are in tension, and I measured it: on my corpus the mean chunk
is 235 tokens, the maximum 263, and 4 of 65 chunks exceed the limit and are
silently truncated --- about two tokens each, so negligible here. But the headroom
is thin, and it only holds because transcript text with interleaved timestamps has
a high characters-per-token ratio. On dense prose the same setting would truncate
most chunks, and the dropped text would be unsearchable with no error raised. The
correct fix is to split on the tokeniser, not on characters.

\item \textbf{\texttt{transcript\_id} is a SHA-256 of the \emph{filename}. What
breaks?}\\
\emph{Answer:} Two different files named \texttt{meeting.txt} collide. Line 141
hashes the name only, \texttt{has\_existing\_index} then reports a hit (line 174),
and the app serves the first file's embeddings while displaying the second file's
name --- confidently wrong answers with no error. It should hash the file
\emph{contents}, which also gives the deduplication the current design was aiming
at. This is a bug I found by auditing rather than by being told.

\item \textbf{Your prompt says ``use ONLY the transcript excerpts''. How do you
know the model obeys?}\\
\emph{Answer:} I do not, and that is the honest gap. I measure retrieval quality
(hit@4 0.83) but not generation faithfulness. Nothing checks the answer against the
context, the UI shows no citations, and there is no refusal test set. To close it I
would build a small labelled set including questions the transcript cannot answer,
and score both refusal behaviour and span attribution --- either with an LLM judge
or by checking that cited timestamps exist in the retrieved chunks.

\item \textbf{If retrieval returns four irrelevant chunks, what happens?}\\
\emph{Answer:} The LLM is called anyway. There is no score threshold and no
relevance filter (line 231), so the prompt instruction is the only thing standing
between the user and a hallucination. The intended fix is exactly the piece I did
not build: a grading node scoring each chunk, with a conditional edge to a
query-rewrite or refuse path --- corrective RAG. That is also the honest answer to
``why isn't this agentic''.

\item \textbf{If the graph is a fixed two-node chain, why use LangGraph at all?}\\
\emph{Answer:} For the checkpointer and streaming, not for routing. Compiling with
\texttt{SqliteSaver} (line 254) gives me thread-scoped conversation state that
survives a browser refresh and a process restart, keyed as
\texttt{\{transcript\_id\}::\{chat\_id\}}, which is what lets a user resume a named
conversation about a transcript they uploaded last week. \texttt{stream\_mode="messages"}
with a node filter lets me stream tokens from \texttt{generate} only. I could have
written both by hand; that is precisely the abstraction-versus-effort trade-off my
third resume bullet is about.

\item \textbf{One module-level SQLite connection with
\texttt{check\_same\_thread=False}, and Streamlit runs a thread per session. Is
that safe?}\\
\emph{Answer:} It is safe for one local user, which is what it was built for, and
it is the first thing that breaks under concurrency. The connection at line 31 is
shared across every session thread; concurrent writes serialise at best and
contend at worst. For a real deployment I would move the checkpointer to Postgres
and give each browser session its own identity, since \texttt{thread\_id} currently
derives from the filename alone.

\item \textbf{Why the map-reduce summariser instead of one call?}\\
\emph{Answer:} A hard constraint, handled deliberately. Groq's free tier caps
\texttt{gpt-oss-120b} at roughly 8000 tokens per minute, and a full transcript in
one prompt exceeds it. So \texttt{\_condense} (lines 115--125) summarises
3000-character windows first and feeds the partial summaries, not the raw
transcript, into the final structured call, with a 1.5\,s pace between calls and
exponential backoff on \texttt{RateLimitError} (lines 95--107). It trades fidelity
for feasibility --- detail is lost in the map step --- and on a paid tier with a
long-context model I would summarise in one pass.

\end{enumerate}

% =====================================================================
\section{Defects found during the audit}\label{sec:defects}
% =====================================================================

Listed in the order they should be fixed.

\begin{enumerate}[leftmargin=*]
  \item \textbf{Five live API keys are committed to git} (\texttt{.env}, present
        in the initial commit, with a GitHub remote configured). OpenAI, Google,
        HuggingFace, Groq and an exchange-rate key. \textbf{Rotate all five},
        add \texttt{.env} to \texttt{.gitignore}, and
        \texttt{git rm -{}-cached .env}. Rotation is what makes the leak harmless;
        deleting the file does not, because the keys remain in history.
  \item \textbf{A Python virtual environment is tracked in git} ---
        \texttt{langchain/} accounts for 18{,}110 of the repository's 18{,}184
        uncommitted changes and is why the editor disables Git features.
  \item \textbf{Filename-based transcript identity} (line 141) serves the wrong
        index on a name collision. Hash contents instead.
  \item \textbf{\texttt{requirement.txt} cannot install the application} --- it is
        missing Streamlit, LangGraph, Chroma, Groq and sentence-transformers, and
        pins nothing.
  \item \textbf{Retrieval ignores conversation history} (line 235), so follow-up
        questions retrieve on pronouns.
  \item \textbf{Character-based chunking against a token-limited encoder}
        truncates 6\% of chunks on this corpus and would truncate most of them on
        prose.
  \item \textbf{\texttt{Chroma} imported from \texttt{langchain\_community}} is
        deprecated in favour of the \texttt{langchain-chroma} package.
\end{enumerate}

% =====================================================================
\section{If you have one more week}\label{sec:nextsteps}
% =====================================================================

Ranked by interview value per hour spent:

\begin{enumerate}[leftmargin=*]
  \item \textbf{Run the harness against the NumPy implementation.} It turns
        bullet 3 from a claim into a result, and it is the single highest-value
        thing left. (Half a day.)
  \item \textbf{Fix chunking and re-measure.} Split on tokens rather than
        characters, try 512 and 256 token chunks, and report how hit@4 moves. You
        would then own a sentence no other candidate has: \emph{``chunk-boundary
        misses were my dominant failure mode, so I changed the splitter and hit@4
        went from 0.83 to $X$.''} (Half a day, and the harness already exists.)
  \item \textbf{Add the grading node.} One conditional edge makes the graph
        genuinely branch, and legitimises far stronger language about it.
        (Half a day.)
  \item \textbf{Deploy it properly on Hugging Face Spaces} and restore the word
        ``deployed'' honestly. (Half a day, plus the persistence decision.)
\end{enumerate}

\vspace{1em}
\noindent
\textbf{One closing note on framing.} The most persuasive thing in this project is
not that it works --- most candidates' projects work. It is that you can state the
measured hit rate, name the dominant failure mode, and explain precisely which
abstraction you took from LangGraph and why. Keep the audited vocabulary:
\emph{served locally}, not deployed; \emph{two-node graph}, not agent;
\emph{context-constrained}, not verified-grounded. Every one of those swaps trades
a word you would lose on for a sentence you can win with.

\end{document}
